\scanpage{35}\mspage{34}
% @old
\deleted{\uline{Contrée finitiste} celle où toute figure}
% @fig

\deleted{\uline{Cont 5}} \inserted{\textcircled{\scriptsize Fig 4}} Soient $F$ et $\unsure{G|C}$ deux figures, telles que pour tt \deleted{$\forall$}$X\in F$, $Y\in G$, $x$ et $Y$ comp.
% @side
% slanted
\margin{On dit alors que $F$ et $G$ sont des \uline{figures compatibles}}
% @union
Alors $F\cup G$ est une figure -- \deleted{\unsure{tels|tel}} (\uuline{figure réunion}).
% @meet

Bien sur $F\cap G$ est une sous-figure commune de $F$ et de $G$.
% @oldcor

\deleted{Corollaire (de Fig 1, Fig 2, Fig 4 -- renforce le cor. précédent) Si $X,Y,\unsure{x'|X'},\unsure{y'|Y'}$ sont tels que $X'\ll X$, $Y'\ll Y$, et $\{X,Y\}$, $\{X',X\}$, $\{Y',Y\}$ comp., alors $X',Y'$ comp.}
% @fin

\uline{Contrée finitiste} Ou les figures sont finies. \deleted{Multistructure finitiste. \unsure{/}} \deleted{Une}
% @cor

\uline{Corollaire} \inserted{(de \uline{Cont} 2, \uline{Cont} 3, \uline{Cont} 5)} Les figures finies \deleted{sont} exactement les parties \uline{finies} de $M$, fermées pour $\leq$, formées d'é\deleted{\unsure{éléments|léments}} \deleted{\unsure{une}} d'é\deleted{\unsure{éléments|léments}} \inserted{\unsure{multisttes|multistés}} qui sont mutuellement compatibles.
% @donc

\unsure{Donc}, \deleted{la structure Contrée finitiste} \deleted{\unsure{une}} \inserted{les figures finies} \unsure{se pvt|se pt} \deleted{\unsure{et co}} connaître quand on connaît seulement $\ll,\unsure{\mathcal R|\mathcal Q}_{M}$. \inserted{Donc une contrée finitiste est connue, par les seules structures $\ll,\mathcal R_{M}$ ;}
% @question
Si on se donne deux telles relations, \inserted{(ordre, et rel. sym. réfle.)} satisfaisant \uline{Cont} 1, à quelles conditions les parties finies $\leq$-fermées \deleted{\unsure{que}} et formées d'éléments mut. compatibles (i.e les \uline{configurations}) satisf\deleted{\unsure{ass|ais}}ont-elles \uline{Cont} 2 à \uline{Cont} 5, ?
